% ============ 第2章 §2.10 Hardy空间的边界值 (III)(原书 pp.149-153) ============
\section{Hardy 空间的边界值 (III)}\label{sec:2.10}

对 $0 < p \leq \infty$,令
\begin{equation*}
H_{*}^{p} = \left\{f(e^{i\theta}) = \lim_{r \to 1}f(re^{i\theta}) \text{ a.e. } \theta;\ f \in H^{p}\right\}.
\end{equation*}
当 $f, g \in H_{*}^{p}$ 且 $f(e^{i\theta}) = g(e^{i\theta})$ a.e. $\theta$ 时,把 $f$ 与 $g$ 视为同一,则由定理~\ref{thm:2.5.3},$H_{*}^{p} \subset L^{p}$.此时,由定理~\ref{thm:2.7.1},$H^{p}$ 与 $H_{*}^{p}$ 可以等同看待.本节证明 $H_{*}^{p}$ 是 $L^{p}$ 的闭子空间,并由此导出 $H^{p}$ 完备.

对 $f \in L^{p}$,令
\begin{equation*}
\|f\|_{p}^{p} = \int_{-\pi}^{\pi}|f(t)|^{p}\frac{\dd t}{2\pi} \quad (p \neq \infty)
\end{equation*}
且 $\|f\|_{\infty} = \esssup_{-\pi \leq t \leq \pi}|f(t)|$.当 $1 \leq p \leq \infty$ 时 $\|f\|_{p}$ 是范数,$L^{p}$ 是 Banach 空间.当 $0 < p < 1$ 时,令 $d(f, g) = \|f - g\|_{p}^{p}$,则是距离,$L^{p}$ 是完备度量空间.

本节证明当 $1 \leq p \leq \infty$ 时 $H_{*}^{p} = \{f \in L^{p};\ \hat{f}(n) = 0 \ (n < 0)\}$.

\begin{lemma}\label{lem:2.10.1}
当 $0 < p \leq \infty$ 时,若 $f \in H^{p}$,则
\begin{equation*}
|f(z)| \leq 2^{1/p}\|f\|_{p}(1 - r)^{-1/p} \quad (r = |z|).
\end{equation*}
\end{lemma}

\begin{proof}
当 $p = 1$ 时,由定理~\ref{thm:2.9.1},
\begin{equation*}
|f(z)| \leq \frac{1}{2\pi}\int_{-\pi}^{\pi}P(r, \theta - t)|f(e^{it})|\dd t \leq \frac{1 + r}{1 - r}\|f\|_{1},
\end{equation*}
故显然.当 $p = \infty$ 时由定理~\ref{thm:2.7.1} 显然,故设 $0 < p < \infty$.由定理~\ref{thm:2.6.4} 可分解 $f = Bg$,$|f(z)| \leq |g(z)|$ 且 $\|f\|_{p} = \|g\|_{p}$.由于 $[g(z)]^{p} \in H^{1}$,由 $p = 1$ 的情形,
\begin{equation*}
|g(z)|^{p} \leq \frac{1 + r}{1 - r}\|g^{p}\|_{1} = \frac{1 + r}{1 - r}\|f\|_{p}^{p},
\end{equation*}
引理得证.
\end{proof}

\begin{theorem}\label{thm:2.10.2}
当 $0 < p \leq \infty$ 时,$H_{*}^{p}$ 是 $L^{p}$ 的闭子空间.当 $p \neq \infty$ 时,$e^{i\theta}$ 的多项式的全体在 $H_{*}^{p}$ 中稠密.
\end{theorem}

\begin{proof}
当 $p \neq \infty$ 时,证明 $e^{i\theta}$ 的多项式的全体在 $H_{*}^{p}$ 中稠密.设 $f \in H_{*}^{p}$ 且 $\varepsilon > 0$,由 $H_{*}^{p}$ 的定义可以把 $f$ 看作属于 $H^{p}$,故由定理~\ref{thm:2.7.3},存在 $0 < \rho < 1$ 使得
\begin{equation*}
\|f_{\rho} - f\|_{p} < \varepsilon.
\end{equation*}
这里 $f_{\rho}(z) = f(\rho z)$.设 $s_{n}(z)$ 为 $f(z)$ 在 $z = 0$ 处 Taylor 展开的第 $n$ 个部分和.此时,由于 $s_{n}(z)$ 在 $|z| = \rho$ 上一致收敛于 $f(z)$,对足够大的 $n$,
\begin{equation*}
\|s_{n\rho} - f_{\rho}\|_{p} < \varepsilon.
\end{equation*}
这里 $s_{n\rho}(z) = s_{n}(\rho z)$.因此,利用当 $0 < p < \infty$ 时
\begin{equation*}
(a + b)^{p} \leq 2^{p}(a^{p} + b^{p}) \quad (a \geq 0, b \geq 0)
\end{equation*}
成立,可以证明 $e^{i\theta}$ 的多项式在 $H_{*}^{p}$ 中稠密.

证明 $H_{*}^{p}$ 是 $L^{p}$ 的闭子空间.设 $f_{n} \in H_{*}^{p}$,$\phi \in L^{p}$ 且 $\|f_{n} - \phi\|_{p} \to 0$($n \to \infty$),则 $\{f_{n}\}$ 在 $L^{p}$ 中有界,故由引理~\ref{lem:2.10.1},对任意 $\rho < 1$,$\{f_{n}(z)\}$ 在 $|z| \leq \rho$ 上一致有界.$\{f_{n}\}$ 在 $|z| \leq \rho$ 上是正规族,故存在其子列 $\{f_{n_{j}}\}$ 使得 $f_{n_{j}}(z)$ 在 $|z| \leq \rho$ 上一致收敛于正则函数 $f(z)$.由一致定理可证 $f$ 与 $\rho$ 无关,且由 $\{f_{n}\}$ 在 $L^{p}$ 中有界知 $f \in H^{p}$.证明 $\phi(\theta) = f(e^{i\theta})$ a.e. $\theta$ 后证明结束.对任意 $\varepsilon > 0$,存在 $N$ 使得当 $n, m \geq N$ 时 $\|f_{n} - f_{m}\|_{p} < \varepsilon$.因此,当 $m \geq N$ 且 $r < 1$ 时,
\begin{align*}
M_{p}(r, f - f_{m}) & = \lim_{j \to \infty}M_{p}(r, f_{n_{j}} - f_{m}) \\
& \leq \lim_{j \to \infty}\|f_{n_{j}} - f_{m}\|_{p} \leq \varepsilon.
\end{align*}
由定理~\ref{thm:2.7.1},当 $m \geq N$ 时 $\|f - f_{m}\|_{p} \leq \varepsilon$.因此 $\|f - f_{n}\|_{p} \to 0$ 且 $\phi(\theta) = f(e^{i\theta})$ a.e. $\theta$.
\end{proof}

下面两个系是定理~\ref{thm:2.7.1} 与上述定理的结果.

\begin{corollary}\label{cor:2.10.3}
当 $1 \leq p \leq \infty$ 时,$H^{p}$ 是 Banach 空间.
\end{corollary}

\begin{corollary}\label{cor:2.10.4}
当 $0 < p < 1$ 时,$H^{p}$ 是完备度量空间.
\end{corollary}

当 $f(z) = \sum_{n=0}^{\infty}a_{n}z^{n} \in H^{2}$ 时,令 $F(e^{i\theta}) = \sum_{n=0}^{\infty}a_{n}e^{in\theta}$,则 $F \in L^{2}$,但更一般地证明 $F(e^{i\theta}) = f(e^{i\theta})$ a.e. $\theta$,即 $\hat{f}(n) = a_{n}$($n \geq 0$).

\begin{theorem}\label{thm:2.10.5}
设 $1 \leq p \leq \infty$.若 $f(z) = \sum_{n=0}^{\infty}a_{n}z^{n} \in H^{p}$,则 $\hat{f}(n) = a_{n}$($n \geq 0$)且 $\hat{f}(n) = 0$($n < 0$).这里 $\hat{f}(n)$ 是 $f(e^{i\theta})$ 的 Fourier 系数.更进一步,
\begin{equation*}
H_{*}^{p} = \left\{\phi \in L^{p};\ \hat{\phi}(n) = 0 \ (n = -1, -2, \ldots)\right\}.
\end{equation*}
\end{theorem}

\begin{proof}
设 $f \in H^{p}$,则当 $0 < r < 1$ 时
\begin{equation*}
\hat{f}_{r}(n) = \frac{1}{2\pi}\int_{-\pi}^{\pi}e^{-int}f(re^{it})\dd t = \begin{cases} r^{n}a_{n} & (n \geq 0) \\ 0 & (n < 0) \end{cases},
\end{equation*}
且 $|\hat{f}_{r}(n) - \hat{f}(n)| \leq \|f_{r} - f\|_{1}$.因此,由定理~\ref{thm:2.7.3},令 $r \to 1$,$\hat{f}(n) = \lim_{r \to 1}\hat{f}_{r}(n)$,定理的前半得证.对后半,由前半只需证明 $H_{*}^{p} \supset \{\phi \in L^{p};\ \hat{\phi}(n) = 0 \ (n < 0)\}$.这由命题~\ref{prop:2.9.3},当 $p = 1$ 时已证明.当 $p \neq 1$ 时也类似.
\end{proof}

在 $H^{p}$ 的边界值的全体 $H_{*}^{p}$ 中,$H_{*}^{1}$ 不包含除常数外的实值函数,由定理~\ref{thm:2.10.5} 是显然的.当 $0 < p < 1$ 时,$H_{*}^{p}$ 中存在除常数外的实值函数.事实上,令 $f(z) = i(1 + z)(1 - z)^{-1}$,则 $f \in \bigcap_{p < 1}H^{p}$,且 $f(e^{i\theta}) = -2\sin\theta|1 - e^{i\theta}|^{-2} \in \bigcap_{p < 1}H_{*}^{p}$.另一方面,$H_{*}^{1/2}$ 不包含除常数外的非负函数 \cite{ch2:1}.下面的命题为了在 2.16 节使用而更一般地叙述 \cite{ch2:2}.

\begin{proposition}\label{prop:2.10.6}
若 $f \in e^{-in\theta}H_{*}^{1/2}$,$n \geq 0$,在 $\partial D$ 上非负,则 $f$ 是次数至多 $n$ 的三角多项式.
\end{proposition}

\begin{proof}
属于 $e^{-in\theta}H_{*}^{1}$ 的非负值函数是次数至多 $n$ 的三角多项式是显然的.设 $f \in e^{-in\theta}H_{*}^{1/2}$ 是非负值函数,为证明命题只需证明 $f \in e^{-in\theta}H_{*}^{1}$.设 $f(e^{i\theta}) = e^{-in\theta}g(e^{i\theta})$ 且 $g \in H_{*}^{1/2}$,由定理~\ref{thm:2.8.3},$g = QG$ 可唯一分解.这里 $Q$ 是内函数,$G$ 是外函数.令 $g_{1} = QG^{1/2}$,$g_{2} = G^{1/2}$,则 $g = g_{1}\cdot g_{2}$,且 $g_{j} \in H_{*}^{1}$($j = 1, 2$).此时 $|g_{1}| = |g_{2}|$.由于 $g_{1}\cdot e^{-in\theta}g_{2} = |g_{1}|^{2}$,故 $e^{-in\theta}g_{2} = \bar{g}_{1}$.因此 $g_{1} + e^{-in\theta}g_{2}$ 与 $i(g_{1} - e^{-in\theta}g_{2})$ 是实值函数.

当 $0 < r < 1$ 时,设 $A_{r} = \{z \in \mathbb{C};\ r < |z| < 1\}$.若 $F$ 在 $A_{r}$ 上正则,
\begin{equation*}
\sup_{r \leq t < 1}\int_{-\pi}^{\pi}|F(te^{i\theta})|\dd\theta < \infty,
\end{equation*}
且 $F$ 在 $\partial D$ 上实值,则 $F$ 可以解析延拓越过 $\partial D$.这个结果超出本书的范围,但众所周知,其证明见 \cite{ch2:3}.

令
\begin{equation*}
F_{1}(z) = g_{1}(z) + \frac{g_{2}(z)}{z^{n}}, \quad F_{2}(z) = i\left(g_{1}(z) - \frac{g_{2}(z)}{z^{n}}\right),
\end{equation*}
则 $F_{1}$ 与 $F_{2}$ 在 $\partial D$ 上正则.因此 $g_{1}$ 在 $\partial D$ 上正则,且由 $g_{1} \in H_{*}^{\infty}$ 得 $g \in H_{*}^{\infty}$.这表明 $f$ 属于 $e^{-in\theta}H_{*}^{1}$.
\end{proof}
